\chapter{Eigenwaarden en eigenvectoren}
\label{ch:eigenwaarden_eigenvectoren}


\section{Wat zijn eigenwaarden en eigenvectoren?}

Stel je een vector $\vec{v}$ voor, en een matrix $\vec{A}$ en een constante $\lambda$.

\[ \vec{A}\vec{v} = \lambda \vec{v} \]

Dus als een matrixvermenigvuldiging van een vector resulteert in een vector die een veelvoud is van de originele vector,
dan noemen we die vector een \textbf{eigenvector} van de matrix $A$, en de constante $\lambda$ noemen we de \textbf{eigenwaarde} van de matrix $A$.

Een matrixvermedigvuldiging heeft dan hetzelfde effect als het schalen van de vector met de factor $\lambda$.

\begin{voorbeeld}[Voorbeeld van eigenwaarden en eigenvectoren]
    
\end{voorbeeld}

